\chapter*{Abstract}
Interindividual variability limits the generalization of motor imagery brain-computer interfaces and makes it difficult to compare architectures evaluated under different protocols. Left- versus right-hand motor imagery classification was evaluated using scalp electroencephalography from OpenBMI, a public dataset of 54 participants, and WBCIC-SHU was used for external validation with 45 common channels. Signals were filtered, cleaned, segmented and normalized before adaptation to four architectures. EEGNet served as the baseline; Conformer and EEGMamba were trained from scratch, whereas LaBraM Base was fully fine-tuned from pretrained weights. Session-wise within-subject evaluation and leave-one-subject-out (LOSO) validation used seeds 2026, 2027 and 2028. An indicator compared seen and unseen participants using the same weights to separate this change from training data volume; external validation additionally examined changes in dataset and protocol. Under LOSO, mean accuracies across 54 participants and three seeds were 0.761 for EEGNet, 0.747 for Conformer, 0.755 for LaBraM and 0.715 for EEGMamba. Optimization regimes were not identical across all seeds. EEGNet had the lowest interparticipant degradation indicator, whereas EEGMamba had the lowest protocol-change indicator. The evaluation included uncertainty, stability and additional criteria for comparative decision-making. The rule, set after the first analysis for the primary dataset and the regimes evaluated, preferred EEGNet because of its lower measured cost among the three architectures equivalent in LOSO accuracy and degradation according to the Bayesian test and Holm-adjusted TOST. This preference does not imply superiority on every axis or beyond the conditions evaluated. These results restrict the interpretation of generalization to the datasets and conditions studied and provide a reproducible basis for model comparison and further research.

\noindent\textbf{Keywords:} brain-computer interface; motor imagery; electroencephalography; domain shift; deep learning; generalization.
